\documentclass[dvipdfmx,aspectratio=169]{beamer}
\usepackage{amsmath}
\usepackage{booktabs}
\usepackage{pxjahyper}
\usepackage{tikz}
\usetikzlibrary{arrows.meta,positioning,backgrounds,decorations.pathreplacing}

\usetheme{Madrid}
\usecolortheme{default}
\setbeamertemplate{navigation symbols}{}

\title[ARC127 A]{AtCoder ARC127 -- A\\\small Leading 1s の解説}
\author{}
\date[ARC127]{AtCoder Regular Contest 127 / 300 点}

\definecolor{ok}{RGB}{0,140,70}
\definecolor{ng}{RGB}{200,30,30}
\definecolor{hi}{RGB}{240,180,0}

\begin{document}

\begin{frame}
  \titlepage
\end{frame}

%====================================================================
\begin{frame}{問題の概要}
  整数 $x$ に対し，\textbf{$f(x) = x$ を十進表記したときの先頭から連続する $1$ の個数}
  と定める．

  \vspace{0.4em}
  \begin{center}
    \begin{tikzpicture}[scale=1.0]
      \foreach \i/\d/\c in {0/1/hi, 1/1/hi, 2/1/hi, 3/0/white, 4/1/white}{
        \draw[thick] (\i*0.85,0) rectangle (\i*0.85+0.75,0.75);
        \fill[\c] (\i*0.85+0.03,0.03) rectangle (\i*0.85+0.72,0.72);
        \draw[thick] (\i*0.85,0) rectangle (\i*0.85+0.75,0.75);
        \node at (\i*0.85+0.375,0.375) {\d};
      }
      \draw[decorate,decoration={brace,amplitude=5pt},ok,thick]
        (0,0.95) -- (2.35,0.95) node[midway,above,ok,font=\small] {連続する $1$ が 3 個};
      \node[anchor=west,font=\small] at (4.5,0.375) {$\Rightarrow f(11101) = 3$};
      \node[anchor=west,font=\scriptsize,ng] at (2.6,-0.35) {ここで途切れる};
    \end{tikzpicture}
  \end{center}

  \vspace{0.4em}
  他の例: $f(1)=1$,\ \ $f(2)=0$,\ \ $f(10)=1$,\ \ $f(11)=2$,\ \ $f(101)=1$

  \vspace{0.5em}
  \textbf{求めるもの}: $f(1)+f(2)+\cdots+f(N)$

  \vspace{0.3em}
  \begin{alertblock}{制約}
    $1 \le N \le 10^{15}$ \quad（高々 16 桁）$\Rightarrow$
    \textbf{1 個ずつ $f$ を計算する $O(N)$ は論外．数え方を変える必要がある．}
  \end{alertblock}
\end{frame}

%====================================================================
\begin{frame}{核心 -- $f$ を $0/1$ 判定の和にバラす}
  $f(x)$ をそのまま計算しようとすると $x$ ごとのループになって詰む．そこで
  \[
    \boxed{\ f(x) = \sum_{k \ge 1} \bigl[\,x \text{ の先頭 } k \text{ 桁がすべて } 1\,\bigr]\ }
  \]
  （$[\cdot]$ は成り立てば $1$，そうでなければ $0$）

  \vspace{0.1em}
  \begin{columns}
    \begin{column}{0.48\textwidth}
      \centering
      {\small
      \begin{tabular}{ccc}
        \toprule
        $k$ & $x=111$ の先頭 $k$ 桁 & 判定\\
        \midrule
        1 & \texttt{1}   & \textcolor{ok}{真}\\
        2 & \texttt{11}  & \textcolor{ok}{真}\\
        3 & \texttt{111} & \textcolor{ok}{真}\\
        4 & （桁が足りない） & \textcolor{ng}{偽}\\
        \bottomrule
      \end{tabular}}

      \vspace{0.3em}
      {\small 真になる $k$ はちょうど $3$ 個 $= f(111)$}

      \vspace{0.2em}
      {\scriptsize ※ 同じ $x$ を複数回数えるのは\textbf{仕様}．\\
      それがちょうど $f(x)$ になる．}
    \end{column}
    \begin{column}{0.48\textwidth}
      和の順序を入れ替えると
      \[
        \text{答え} = \sum_{k \ge 1} \#\{\,x \le N : \text{先頭 } k \text{ 桁が全部 } 1\,\}
      \]

      \vspace{0.3em}
      \begin{block}{何が起きたか}
        \textbf{$x$ のループ（$10^{15}$ 回）が\\
        $k$ のループ（高々 16 回）に化けた．}
      \end{block}
    \end{column}
  \end{columns}
\end{frame}

%====================================================================
\begin{frame}{先頭が固定された数の集合 $=$ 区間の集まり}
  $\mathrm{head} = \underbrace{11\cdots1}_{k \text{ 個}}$ とおく．
  「先頭 $k$ 桁が $1$ で，後ろに $i$ 桁続く数」は，$L = \mathrm{head}\times 10^{i}$ として
  \[
    \boxed{\ [\,L,\ L + 10^{i} - 1\,] \quad \text{のちょうど } 10^{i} \text{ 個}\ }
  \]

  \vspace{0.2em}
  \begin{center}
    \begin{tikzpicture}[scale=1.0]
      \draw[-{Latex},thick] (-0.3,0) -- (11.0,0) node[right]{};
      % i=0
      \fill[ok] (0.6,-0.08) rectangle (0.75,0.08);
      \node[font=\scriptsize,anchor=south] at (0.7,0.15) {111};
      \node[font=\scriptsize,anchor=north,ok] at (0.7,-0.2) {$i{=}0$: 1 個};
      % i=1
      \draw[ok,line width=6pt] (2.6,0) -- (3.5,0);
      \node[font=\scriptsize,anchor=south] at (3.05,0.2) {1110 \textasciitilde\ 1119};
      \node[font=\scriptsize,anchor=north,ok] at (3.05,-0.2) {$i{=}1$: 10 個};
      % i=2
      \draw[ok,line width=6pt] (5.6,0) -- (8.0,0);
      \node[font=\scriptsize,anchor=south] at (6.8,0.2) {11100 \textasciitilde\ 11199};
      \node[font=\scriptsize,anchor=north,ok] at (6.8,-0.2) {$i{=}2$: 100 個};
      \node[font=\scriptsize,anchor=west] at (9.2,0) {$\cdots$};
      \node[font=\small,anchor=west] at (-0.3,0.9) {例: $k=3$（$\mathrm{head}=111$）};
    \end{tikzpicture}
  \end{center}

  \vspace{0.6em}
  桁数が違うので，$i$ が違う区間どうしは\textbf{重ならない}．
  よって各 $k$ について「区間の個数を足すだけ」でよい．
\end{frame}

%====================================================================
\begin{frame}{$N$ で打ち切る -- 3 つのケース}
  区間 $[L,\ R-1]$（$R = L + 10^{i}$）に対し，$N$ 以下の部分だけを数える．

  \vspace{0.4em}
  \begin{center}
    \begin{tikzpicture}[scale=1.0]
      % ケース1
      \draw[-{Latex}] (0,2.0) -- (9.5,2.0);
      \draw[ok,line width=6pt] (1.0,2.0) -- (3.2,2.0);
      \draw[ng,thick,dashed] (7.0,1.6) -- (7.0,2.4) node[above,ng,font=\scriptsize]{$N$};
      \node[font=\scriptsize,anchor=west] at (9.6,2.0) {};
      \node[font=\scriptsize,anchor=east] at (-0.1,2.0) {(A)};
      \node[font=\scriptsize,anchor=south] at (2.1,2.15) {丸ごと入る};
      \node[font=\scriptsize,anchor=west,ok] at (3.4,2.0) {$+\,10^{i}$ 個 \ (\texttt{R <= n+1})};
      % ケース2
      \draw[-{Latex}] (0,1.0) -- (9.5,1.0);
      \draw[ok,line width=6pt] (5.0,1.0) -- (7.0,1.0);
      \draw[gray!60,line width=6pt] (7.0,1.0) -- (8.6,1.0);
      \draw[ng,thick,dashed] (7.0,0.6) -- (7.0,1.4) node[above,ng,font=\scriptsize]{$N$};
      \node[font=\scriptsize,anchor=east] at (-0.1,1.0) {(B)};
      \node[font=\scriptsize,anchor=west,ok] at (8.8,1.0) {$+\,(N-L+1)$};
      \node[font=\scriptsize,anchor=south] at (6.0,1.15) {途中で $N$ に到達};
      % ケース3
      \draw[-{Latex}] (0,0.0) -- (9.5,0.0);
      \draw[gray!60,line width=6pt] (7.8,0.0) -- (9.2,0.0);
      \draw[ng,thick,dashed] (7.0,-0.4) -- (7.0,0.4) node[above,ng,font=\scriptsize]{$N$};
      \node[font=\scriptsize,anchor=east] at (-0.1,0.0) {(C)};
      \node[font=\scriptsize,anchor=south] at (8.5,0.15) {$L > N$};
      \node[font=\scriptsize,anchor=west,ng] at (9.4,0.0) {$+\,0$};
    \end{tikzpicture}
  \end{center}

  \vspace{0.3em}
  (B)(C) はどちらも\textbf{以降の $i$ でも $L$ が 10 倍ずつ大きくなる}ので，
  そこで \texttt{break} してよい．
\end{frame}

%====================================================================
\begin{frame}[fragile]{擬似コード}
\begin{verbatim}
ans = 0 ; head = 0
for k in 1..16:                   # N <= 10^15 は高々 16 桁
    head = head * 10 + 1          # 1, 11, 111, ...
    L = head
    for i in 0..15:
        if L > N: break           # (C) この先はすべて範囲外
        R = L + 10^i              # 区間は [L, R-1]
        if R <= N + 1:
            ans += R - L          # (A) 丸ごと入る = 10^i 個
        else:
            ans += N - L + 1      # (B) 途中まで
            break
        L *= 10
\end{verbatim}

  \vspace{0.2em}
  \textbf{計算量} $O(D^2) \approx 16\times 16 = 256$ 回（$D$ = 桁数）．
  \textbf{$N$ の大きさに依存しない．}
\end{frame}

%====================================================================
\begin{frame}{サンプル 2（$N=120$，答え $44$）を手で追う}
  \centering
  {\small
  \begin{tabular}{ccccclr}
    \toprule
    $k$ & head & $i$ & $L$ & 区間 $[L, R-1]$ & 判定 & 加算\\
    \midrule
    1 & 1 & 0 & 1 & $[1,1]$ & (A) 丸ごと & $+1$\\
    1 & 1 & 1 & 10 & $[10,19]$ & (A) 丸ごと & $+10$\\
    1 & 1 & 2 & 100 & $[100,199]$ & (B) $120$ で打ち切り & $+21$\\
    \midrule
    2 & 11 & 0 & 11 & $[11,11]$ & (A) 丸ごと & $+1$\\
    2 & 11 & 1 & 110 & $[110,119]$ & (A) 丸ごと & $+10$\\
    2 & 11 & 2 & 1100 & --- & (C) $L>120$ & break\\
    \midrule
    3 & 111 & 0 & 111 & $[111,111]$ & (A) 丸ごと & $+1$\\
    3 & 111 & 1 & 1110 & --- & (C) $L>120$ & break\\
    \midrule
    4 & 1111 & 0 & 1111 & --- & (C) $L>120$ & break\\
    \bottomrule
  \end{tabular}}

  \vspace{0.6em}
  \raggedright
  小計は $k=1$: $\mathbf{32}$，$k=2$: $\mathbf{11}$，$k=3$: $\mathbf{1}$
  $\Rightarrow$ 合計 $32+11+1 = \mathbf{44}$ \ \textcolor{ok}{\checkmark}

  {\small それぞれ「先頭が $1$ の数の個数」「先頭が $11$ の数の個数」「先頭が $111$ の数の個数」．}
\end{frame}

%====================================================================
\begin{frame}[fragile]{落とし穴 -- \texttt{if L > n: break} を省くと静かに壊れる}
  \begin{columns}
    \begin{column}{0.52\textwidth}
      $L > N$ のとき，(B) の枝が実行されると
      \[
        N - L + 1 \le 0
      \]
      すなわち\textbf{負の数を足してしまう．}

      \vspace{0.6em}
      引き算で個数を出すコードは，範囲外ガードを忘れると
      \textbf{例外も出さずに答えだけが壊れる}．
      個数を出す式を書いたら
      「\textbf{この式が負になる入力はあるか}」を必ず確認する．
    \end{column}
    \begin{column}{0.46\textwidth}
      \begin{block}{実測（ガードを外した版）}
        \begin{itemize}
          \item $N=1$ の出力: \\ \texttt{-1234567901234542}（正解 $1$）
          \item $1 \le N \le 500$ の\\ \textbf{500 通りすべてが不正解}
        \end{itemize}
      \end{block}
    \end{column}
  \end{columns}

  \vspace{0.5em}
  その他の注意点:
  \begin{itemize}
    \item \verb|R <= n + 1| は \verb|R - 1 <= n| の言い換え（区間の右端が $N$ 以下）．
    \item \textbf{Python の int は多倍長なので溢れない．C++ 移植時は \texttt{long long}}
      （$N$ も $L$ も $10^{15}$ 規模で，\texttt{int} では即オーバーフロー）．
    \item ループ上限 16 は $N \le 10^{15}$ が 16 桁だから．制約が $10^{18}$ なら 19 に増やす．
  \end{itemize}
\end{frame}

%====================================================================
\begin{frame}{まとめ}
  \begin{enumerate}
    \item \textbf{$\sum f(x)$ が重いときは $f$ を $0/1$ 判定の和に分解し，
      和の順序を入れ替える}（寄与の分解）．
      \[
        \sum_x f(x) = \sum_x \sum_k [\,\text{条件}_k(x)\,] = \sum_k \#\{x : \text{条件}_k\}
      \]
      $x$ のループが「条件のループ」に化けて，$N$ に依存しない計算量になる．
    \item \textbf{接頭辞が固定された数の集合 $=$ 区間の和．}
      $[\mathrm{head}\cdot 10^i,\ \mathrm{head}\cdot 10^i + 10^i - 1]$ を $i$ について足す．
      桁 DP を持ち出すのは過剰．
    \item \textbf{区間は半開 $[L, R)$ で持つ}と個数が $R-L$ になり，off-by-one が減る．
    \item \textbf{「範囲外なら break」のガードを省かない．}
      引き算で個数を出すコードは負の数を足して静かに壊れる．
  \end{enumerate}

  \vspace{0.4em}
  \textbf{同系統の問題}: ABC154 E -- Almost Everywhere Zero，
  ABC029 D -- 1（桁ごとの寄与），ABC192 D -- Base n
\end{frame}

\end{document}
